% Vienna Basin geothermal --- physics retreat project notes
% Working draft; references collected via web search, May 2026
% Not yet circulated --- to be reviewed before sharing with the group

\section*{Geothermal Heat for Vienna: Seven Research Threads}

\subsection*{Why this scope}

Our preceding analysis converged on geothermal \emph{heat} (not electricity)
for Vienna's district-heating network as the most defensible target for a
physics retreat investigation. At Vienna Basin reservoir temperatures
($\approx 85$--$150\,^\circ$C), the exergy argument strongly favours direct
heat use, and Vienna's $\approx 6$~TWh/yr district-heat demand can in principle
be covered by $\sim 40$ doublets of the kind already being drilled in
Aspern--Donaustadt.

The complication is that this top-level story is already being executed by
Wien Energie and OMV through their joint venture \texttt{deeep Tiefengeothermie
GmbH} (established 2023): three exploratory wells were drilled by mid-2025 at
the Aspern site, surface facilities are being built from 2026, and a $20$~MW$_{\mathrm{th}}$
plant is scheduled to come online in 2028, with up to $\sim 120$~MW$_{\mathrm{th}}$
of capacity planned across Donaustadt and Simmering by 2030 [1, 2]. Simply
reproducing this engineering case study is therefore an exercise, not research.

The seven threads below were chosen specifically because each addresses a
question the utility roadmap does \emph{not} fully cover, while remaining
tractable within a retreat-then-followup timescale. Threads 1--4 are the
original set drafted earlier; threads 5--7 were added in a second pass and
deliberately reach into adjacent physics (systems-level comparison,
working-fluid thermodynamics, multi-decade reservoir behaviour).

% ---------------------------------------------------------------------------
\subsection*{Thread 1 --- Repurposing OMV legacy oil and gas wells}

\paragraph{State of the art.}
Repurposing of legacy hydrocarbon wells for geothermal use has become a
distinct research subfield in the last five years. A 2025 review article
in \emph{Applied Energy} catalogues the four main technical concepts
[11]: (i) direct conversion to an open-loop production
well after re-completion; (ii) coaxial closed-loop single-well borehole heat
exchangers (BHE); (iii) CO$_2$-based thermosiphon designs descended from
Brown (2000); and (iv) use as observation/monitoring wells in a larger
field. Closed-loop coaxial BHEs are the most actively studied because they
sidestep the formation-water risk that complicates open-loop reuse.
Reported single-well thermal outputs are modest: $\sim 600$~kW for a
4~km coaxial BHE under optimal conditions in a numerical study
[12], up to $\sim 2$~MW$_\mathrm{th}$ at early times in the
Mol field case study in Belgium [13]. North American activity is
concentrated in jurisdictions with large orphaned-well inventories
(Alberta, Oklahoma, Colorado), with Gradient Geothermal among the first
commercial movers in the US.

\paragraph{Vienna Basin specifics.}
The Vienna Basin has hundreds of legacy hydrocarbon wells drilled by OMV
and its predecessor RAG since 1949, concentrated in the Matzen,
Sch\"onkirchen, Aderklaa and Bockflie{\ss} fields [5].
Crucially, OMV has \emph{already} initiated a Vienna-Basin pilot:
the Aderklaa-96 well, originally a $\sim 2900$~m gas well drilled in
the 1970s, was re-opened in 2022--2023 for geothermal testing
[3]. Production and injection testing yielded a bottom-hole
temperature of $\approx 102\,^\circ$C in the Hauptdolomit basement and a
reinjection rate up to $\sim 40$~m$^3$/h [4]. This is a live
case study a retreat can build around --- the well exists, the data are
partly public, the depth and temperatures are representative of a much
larger inventory.

\paragraph{Open questions for the retreat.}
\begin{itemize}
\item For Aderklaa-96 specifically: what \emph{steady-state} thermal output
is achievable in coaxial closed-loop mode versus the measured open-loop
performance? Build a 1-D BHE heat-transfer model with the published
geometry and Hauptdolomit thermal properties.
\item Statistical question: given the OMV well inventory in the basin
(depth, temperature, casing condition), what fraction is plausibly
repurposable? Even rough screening criteria would be useful.
\item Wellbore-integrity physics: how does 40--60 years of cement aging,
casing corrosion, and possible $\mathrm{H_2S}$ exposure constrain
re-completion options at OMV-typical depths (1.5--3~km)?
\item CO$_2$ as working fluid: in the closed-loop coaxial geometry,
how much does the supercritical CO$_2$ thermosiphon effect change
thermal output versus water? Vienna Basin temperatures sit close to
the supercritical CO$_2$ regime --- the comparison is non-trivial.
\end{itemize}

% ---------------------------------------------------------------------------
\subsection*{Thread 2 --- Induced seismicity at Vienna Basin conditions}

\paragraph{Why this is not the same problem as Basel or Pohang.}
The Basel 2006 ($M_L\,3.4$) and Pohang 2017 ($M_w\,5.5$) cautionary cases
were Enhanced Geothermal System (EGS) projects in crystalline basement,
involving high-pressure stimulation that reactivated previously unknown
faults. A Vienna Basin hydrothermal project at $\sim 3$~km in sedimentary
host rock, with modest injection over-pressure, is a fundamentally
different operating regime. The relevant analog is the
\emph{Bavarian Molasse Basin} fleet of low-enthalpy hydrothermal plants
(Unterhaching, Pullach, Sauerlach, Geretsried, Poing) ---
geologically and operationally far closer to what Wien Energie plans.

\paragraph{What the Molasse Basin tells us.}
The Poing plant near Munich produced a sequence of $M \approx 2.1$ induced
events about five years after circulation began, with microseismicity
spatially correlated with the injection well
[9]. Numerical ground-motion simulations indicate that
even the largest of these events did not exceed the German
$\mathrm{PGV} = 5$~mm/s damage threshold. The Greater Munich
\emph{INSIDE} research project (BMWi-funded) has since taken on the
broader question of how to monitor and forecast seismicity in a basin
where multiple plants share the same fault system [10].
The headline result: felt induced events are possible from
hydrothermal-scale injection, but stay well below the EGS regime in
both magnitude and damage potential.

\paragraph{Vienna Basin specifics.}
The Vienna Basin is a Miocene pull-apart basin transected by the
left-lateral Vienna Basin Transfer Fault (VBTF). Historical seismicity
along the VBTF reaches $I_{\max,\mathrm{obs}}/M_{\max,\mathrm{obs}} = 8/5.2$,
with major events including the 4th-century AD Carnuntum earthquake, the
2000 Ebreichsdorf $M\,4.8$, the 2013 Ebreichsdorf series ($M\,4.3$),
and a 2024 April $M\approx 3$ event likely on a normal splay fault
of the VBTF [6, 7, 8].
The Hintersberger \& Decker (2018) palaeoseismological study of the
Markgrafneusiedl Fault is particularly relevant: it argues that splay
faults of the VBTF are active despite the \emph{absence} of historical
events on them, and that $M \geq 7.0$ earthquakes should be considered
in hazard assessment for the Vienna metropolitan area [6].
This is a strong reason to take the seismicity question seriously even
under "modest" hydrothermal operating conditions.

\paragraph{Open questions for the retreat.}
\begin{itemize}
\item Build a simple Coulomb-stress / pore-pressure-diffusion model
(Shapiro-style) for a 3~km doublet operating at Aspern, and compare
the resulting failure-pressure perturbation against the static stress
budget of nearby VBTF splay faults.
\item Quantitative comparison Basel/Pohang vs.\ Molasse vs.\ projected Vienna
operating envelope --- a clean physics figure (injection pressure,
flow rate, host rock $k_{\mathrm{perm}}$) suitable for a retreat talk.
\item Did the decades of OMV oil and gas injection in the Vienna Basin
produce any documented induced seismicity? This is a literature
question worth checking systematically against the GeoSphere Austria
catalog --- a "no" answer is itself a strong physical constraint.
\item Design proposal for a microseismic traffic-light system specific
to the Vienna Basin geology, with publicly justifiable thresholds.
\end{itemize}

% ---------------------------------------------------------------------------
\subsection*{Thread 3 --- Lithium (and trace element) co-production}

\paragraph{State of the art.}
Co-production of critical raw materials from geothermal brine is a fast
moving field driven by the EU Critical Raw Materials Act (2024) and
recovering battery demand. The European benchmark is Vulcan Energy
Resources in the Upper Rhine Graben (URG): brines at $\sim 150$--$200$~mg/L
Li, an aluminium-sorbent direct-lithium-extraction (DLE) process, and
a EUR~2.2~billion financing package closed in December 2025.
Commercial production of $\sim 24{,}000$~t/yr lithium-hydroxide-monohydrate
(LHM) is targeted for 2028 [14, 15].

The DLE technology landscape is consolidating around three families:
aluminium-based sorbents (Eramet, Vulcan; $\sim 90\%$ recovery, operates
at native brine temperature), manganese-oxide spinel sorbents (LMO),
and lithium-ion sieve composites [16, 17]. Sorbent
lifetime remains the dominant operating-cost driver. Lithium carbonate
prices recovered sharply in early 2026 (from $\sim\$13$k/t in Dec 2025
to $\sim\$26$k/t in late January, with Chinese prices around
CNY~175{,}000/t in May 2026), and forecasts shift the market from
surplus to potential supply deficit through 2026
[18].

\paragraph{Vienna Basin specifics --- the genuine research gap.}
This thread contains the most candid research opening of the four:
\emph{no published lithium concentration data for Vienna Basin formation
waters was found in the searches conducted for this draft}. Brine
chemistry studies for the basin exist in the petroleum-geology
literature (TDS, salinity, gas composition) and the general theory of
lithium enrichment in sedimentary formation waters is well developed
[19], but Li-specific assays for OMV produced
water from the Matzen / Sch\"onkirchen / Aderklaa fields appear to
be either unpublished or buried in operator archives. The URG Li
concentrations are the product of a very specific
felsic-basement / Permo-Triassic-aquifer setting; the Pannonian-edge
Vienna Basin has a different sediment provenance and may well have
much lower Li. But until someone actually \emph{measures} this, the
question is open.

\paragraph{Open questions for the retreat.}
\begin{itemize}
\item Literature triage: produce a systematic review of all Vienna Basin
formation-water chemistry studies (GBA / GeoSphere Austria reports,
OMV production geology theses from TU Wien and Montanuniversit\"at
Leoben) and tabulate any Li, B, Cs, Rb measurements that exist.
\item Order-of-magnitude estimate: even at URG-like Li concentrations
($\sim 175$~mg/L) and Aderklaa-96-like flow rates ($\sim 40$~m$^3$/h),
what would the annual LCE output be? Is it economic at 2026 prices?
\item Geochemical-equilibrium reasoning: based on Vienna Basin host
lithologies (carbonates, Miocene clastics) and reservoir temperatures,
what Li concentration range is \emph{plausible} a priori?
\item Helium and noble gases: any indication of He enrichment in the
deep Vienna Basin that would change the co-product economics?
\end{itemize}

The retreat could plausibly produce an actual unpublished result here
(an order-of-magnitude bound on Vienna Basin Li potential), rather than
just a literature review.

% ---------------------------------------------------------------------------
\subsection*{Thread 4 --- Seasonal storage coupling (HT-ATES)}

\paragraph{State of the art.}
The seasonality of Vienna's heat demand (peak $\sim 3$~GW winter, $\sim 0.3$~GW
summer) means that constant-output geothermal is over-built for winter or
under-utilised in summer unless coupled to seasonal storage.
High-Temperature Aquifer Thermal Energy Storage (HT-ATES), at
$\geq 60\,^\circ$C and up to $\sim 90\,^\circ$C, is the technology most
likely to integrate cleanly with deep geothermal because both rely on
the same hydrogeological substrate [20]. The EU
\emph{HEATSTORE} project has supported pilot installations in
Switzerland (Bern Forsthaus, Geneva), the Netherlands and Iceland;
TU Delft operates a campus-scale HT-ATES research demonstrator;
the Burgwedel (Germany) techno-economic case study from 2024 is a
useful benchmark for HT-ATES economics in conditions resembling
Vienna Basin sediments [21].

The dominant physics issues are: (i) round-trip efficiency, typically
60--80\% and declining with storage temperature; (ii) buoyancy-driven
tilting of the thermal plume in HT regime; (iii) thermo-hydro-mechanical
coupling and small ground-surface uplift [22];
(iv) hydrochemical scaling/clogging at higher temperatures.

\paragraph{Vienna Basin specifics.}
Wien Energie has an ongoing \emph{ATES Vienna} investigation under the
Green Energy Lab umbrella, examining thermal-water reservoirs at
600--3000~m depth for seasonal district-heat storage --- explicitly
framed as a prerequisite for full DH decarbonisation [23].
The project is in an early/exploratory phase; quantitative results
on aquifer suitability, achievable storage temperatures, and
round-trip efficiency for the specific Vienna Basin aquifers are not
yet in the literature. There is therefore both clear utility interest
in the question and room for an independent academic contribution.

\paragraph{Open questions for the retreat.}
\begin{itemize}
\item \textbf{Headline optimisation problem.} Given Vienna's annual
heat-demand curve (well known) and a parameterised HT-ATES (storage
volume $V$, storage temperature $T_{\mathrm{st}}$, round-trip
efficiency $\eta(T_{\mathrm{st}})$), what combination of geothermal
doublet count $N$ and storage volume $V$ minimises the levelised
cost of heat (LCOH) for a given decarbonisation target? This is a
clean, retreat-sized optimisation problem with publishable potential.
\item Heat-loss scaling: confirm that the surface-to-volume scaling
for an HT-ATES at Vienna-relevant scales gives the round-trip
efficiencies cited in HEATSTORE, or identify the regime where it
does not.
\item Thermal-plume migration and buoyancy: under what aquifer
parameters does buoyancy-driven tilting become a first-order loss
mechanism? A 2-D simulation with realistic Vienna Basin permeability
contrasts would be illuminating.
\item Alternative storage technologies for comparison: pit thermal
energy storage (PTES, cf.\ Danish Vojens at $\sim 200{,}000$~m$^3$)
and large-scale heat pumps with thermochemical or hydrogen storage
--- is HT-ATES actually the right answer for Vienna?
\end{itemize}

% ---------------------------------------------------------------------------
\subsection*{Thread 5 --- Geothermal heat vs.\ large heat pumps + green electricity}

\paragraph{Why this is the question Wien Energie can't answer for itself.}
Wien Energie is already operating Europe's most powerful large-scale heat
pump --- a 55~MW$_{\mathrm{th}}$ wastewater-source plant at the Ebswien
treatment site near Simmering, commissioned 2023, which lifts wastewater
from $\sim 11$--$12\,^\circ$C in winter (24--$25\,^\circ$C summer) to
$\sim 90\,^\circ$C feed temperature. The plant currently supplies
$\sim 56{,}000$ households; a second stage doubling capacity to
$110$~MW$_{\mathrm{th}}$ (target $\sim 112{,}000$ households) is scheduled
for completion in 2027 [24, 25]. Capital cost of stage 1 was
$\sim$EUR~70~million, i.e.\ $\sim$EUR~1.3~million per MW$_{\mathrm{th}}$
installed. \emph{The same utility is therefore simultaneously building
deep geothermal and large heat pumps to decarbonise the same district-heat
network.} The implicit question --- in what mix --- is not addressed in any
public-facing document we have located, and the academic literature on
direct LCOH/LCOA comparisons for large-DH applications is surprisingly
thin [26].

\paragraph{The physics-level comparison.}
The two technologies sit at very different points of the
exergy/capex/opex Pareto front:
\begin{itemize}
\item \textbf{Deep geothermal doublet:} CapEx $\sim$EUR~20--35~M per
$\sim 18$~MW$_{\mathrm{th}}$ doublet (i.e.\ $\sim$EUR~1.4~M/MW), capacity
factor $\sim 85$--$90\%$, almost zero operating energy input, $\sim 30$~yr
nominal lifetime with depletion physics (see thread 7), small marginal
cost per MWh.
\item \textbf{Large heat pump:} CapEx $\sim$EUR~1.3~M/MW based on Vienna's
own numbers [24], COP $\sim 3$--$3.5$ for wastewater-source $\to 90\,^\circ$C
delivery, so electricity input $\sim 0.3$~MWh per MWh of heat delivered.
At an Austrian grid carbon intensity of $\sim 100$~g~CO$_2$/kWh
(hydropower-dominated) this is essentially low-carbon today; whether it
stays low-carbon depends entirely on the marginal-electricity assumption
during winter peak.
\end{itemize}
The crossover is dominated not by capex (similar per-MW) but by
\emph{electricity-cost expectations} and \emph{carbon-intensity of
marginal electricity}. Both are policy- and weather-dependent in ways
that make sensitivity analysis essential.

\paragraph{Open questions for the retreat.}
\begin{itemize}
\item Least-cost optimisation over Vienna's annual DH load curve, with
both technologies plus biomass + existing waste-incineration (Spittelau,
already in the network) + HT-ATES (thread 4) competing for share.
What does the optimum mix actually look like?
\item Sensitivity to electricity price (EUR~50--200/MWh) and marginal
carbon intensity (50--500~g~CO$_2$/kWh) --- which of these flips the
ordering?
\item Why is Wien Energie hedging? Is it a portfolio-risk argument
(no single technology covers all of Vienna's $\sim 3$~GW winter peak),
or a thermal-temperature argument (heat pumps marginal for 90~$^\circ$C
output, geothermal at native temperature)?
\item Land use, water use, public acceptance --- secondary axes that
the retreat could quantify rather than just qualitatively note.
\end{itemize}

% ---------------------------------------------------------------------------
\subsection*{Thread 6 --- CO$_2$ (and supercritical) working fluids in closed-loop geothermal}

\paragraph{Why this just became urgent.}
On 4 December 2025, Eavor Technologies delivered its first grid
electricity from the Geretsried (Bavaria) closed-loop multilateral
system --- the first commercial demonstration of an Eavor-Loop at scale
[27, 28]. The Geretsried plant comprises two vertical wells with six
horizontal laterals sidetracked from each, the laterals connected
toe-to-toe underground using active magnetic ranging, forming a
$\sim 12$-lateral subsurface radiator designed for $64$~MW$_{\mathrm{th}}$
and $8.2$~MW$_{\mathrm{e}}$. Crucially, the thermosiphon effect was
established within $\sim 30$ minutes of start-up and the system now
runs pump-free [28]. This is the first time the closed-loop thermosiphon
geometry --- proposed for working fluids including supercritical
CO$_2$ since Brown (2000) --- has been operationally demonstrated as
a power source.

\paragraph{The CO$_2$-versus-water question.}
Eavor uses water; the deeper question is whether supercritical CO$_2$
($T_c = 31.1\,^\circ$C, $p_c = 7.4$~MPa) is a better working fluid at
Vienna-Basin-relevant temperatures. The physics arguments for sCO$_2$:
\begin{itemize}
\item Steeper density gradient with temperature than water in the
relevant regime, hence stronger thermosiphon driving force and lower
parasitic pumping load;
\item Lower viscosity, higher mass flow per unit pressure drop;
\item Potential for higher conversion efficiency in direct-cycle
turbomachinery (Brayton cycle on sCO$_2$);
\item In the open-loop variant (CO$_2$ Plume Geothermal, CPG), permanent
sequestration of CO$_2$ in the host sedimentary basin as a co-benefit [29].
\end{itemize}
The physics arguments against: corrosivity (limits material choices),
high operating pressure ($> 10$~MPa typical), and --- in CPG mode ---
displacement of native brine into the wellbore and potential reactivity
with carbonate host rocks (e.g.\ the Aderklaa Hauptdolomit). For closed
loops, leakage of CO$_2$ to atmosphere defeats the climate purpose.

\paragraph{Vienna Basin specifics.}
No published Vienna Basin sCO$_2$ work has been located in the searches
conducted here. The most concrete physics-flavoured connection is the
overlap with thread 1: the Aderklaa-96 retrofit is, in principle, a
candidate for either a coaxial-water BHE or a single-well sCO$_2$
thermosiphon trial. The Hauptdolomit carbonate host raises the CPG
reactivity question concretely. Multi-lateral Eavor-style geometries are
probably not retrofittable into legacy vertical OMV wells without major
sidetracking work, but the toe-to-toe lateral connection technology is
now demonstrated at Geretsried and could in principle be applied
greenfield in the Vienna Basin.

\paragraph{Open questions for the retreat.}
\begin{itemize}
\item Comparative thermodynamic model: water vs.\ sCO$_2$ in a coaxial
closed-loop at Aderklaa-96 conditions ($T_{\mathrm{prod}}\sim 100$\,$^\circ$C,
$\sim 3$~km depth). Quantify the thermosiphon strength and net power
delivered for each.
\item Brayton-cycle conversion efficiency for sCO$_2$ at the modest
temperatures of the Vienna Basin --- is the efficiency advantage real
at $T_{\mathrm{prod}}<150\,^\circ$C, or does it only show up at the
$\geq 200\,^\circ$C Eavor-style EGS regime?
\item CPG variant: rock--CO$_2$--brine geochemistry in the Aderklaa
Conglomerate and Hauptdolomit --- back-of-envelope reactivity calculation
to identify show-stoppers.
\item Material compatibility of OMV legacy casings (carbon steel, various
vintages) with sCO$_2$ at $\sim 100\,^\circ$C and $\sim 10$~MPa ---
literature review with quantitative degradation rates.
\end{itemize}

% ---------------------------------------------------------------------------
\subsection*{Thread 7 --- Long-term reservoir depletion and multi-doublet interference}

\paragraph{Why this matters for the 40-doublet vision.}
The $\sim 40$ doublets implied by our heat-balance estimate are not
hydraulically independent. Each injection well establishes a thermal plume
in the aquifer that grows with time; eventually the cold front reaches
the production well (``thermal breakthrough''), the produced temperature
declines, and the doublet must be re-sited, re-pointed, or sidetracked.
For Vienna's $\sim 30$-year deployment horizon, the doublet count is
therefore not the only design variable --- the time-evolution of the
reservoir is.

\paragraph{The empirical baseline: Dogger aquifer, Paris Basin.}
The single best dataset for low-enthalpy doublet behaviour over a
human-relevant timescale is the Dogger aquifer in the \^Ile-de-France
($\sim 1500$~m depth, $\sim 60$--$80\,^\circ$C), where district-heating
doublets have been operating continuously since the 1970s. The Lopez
\emph{et al.}\ (2010) retrospective on 40 years of operation [30]
documents the headline result: \emph{none of the active doublets has
shown a statistically significant production-temperature decline} to
date, although modelling clearly shows that the natural heat flux at
the doublet scale is insufficient to maintain current withdrawal rates
indefinitely. A 2024 reanalysis of the Dogger production data with
modern thermo-hydraulic codes [31] confirms that the operationally
relevant decline regime begins at $\sim 50$--$70$~yr for typical
well spacings of 1--1.5~km.

\paragraph{The physics one layer down.}
At doublet scale the thermal breakthrough time scales (roughly) as
$t_{\mathrm{TB}} \propto \frac{(\rho c)_{\mathrm{rock}}\, h\, L^2}{(\rho c)_{\mathrm{water}}\, Q}$
for spacing $L$, aquifer thickness $h$, and flow rate $Q$. This is a
simple analytical model the retreat could parameterise for the
Aderklaa Conglomerate and the Hauptdolomit. The interesting complications
are: (i) facies heterogeneity --- the conglomerate is not a uniform slab;
(ii) cooling-induced poroelastic stress changes large enough to reactivate
distant faults [32], which couples this thread directly to thread 2; and
(iii) cumulative thermal interference between $\sim 40$ doublets across
the Vienna Basin once deployment is mature. In the CO$_2$ Plume Geothermal
variant (thread 6), heat depletion follows different physics again [29],
making CPG potentially more or less sustainable than water doublets
depending on parameters.

\paragraph{Open questions for the retreat.}
\begin{itemize}
\item Analytical thermal-breakthrough estimate for an Aspern doublet with
Aderklaa Conglomerate properties (porosity $\sim 20$\,\%, thickness
$\sim 100$\,m, $L \sim 1.5$\,km) and the Aderklaa-96 reinjection rate
$\sim 40$~m$^3$/h. How does the predicted breakthrough time compare with
the Dogger observation that 50+ years have passed without measurable
decline?
\item Multi-doublet interference: at what doublet density does the
collective thermal field interaction become a first-order constraint on
the maximum sustainable extraction rate per km$^2$ of basin? A 2-D
heat-diffusion simulation with a realistic Vienna Basin geometry would
be illuminating.
\item Coupling to thread 2: estimate the cooling-induced
poroelastic Coulomb stress change at the depth of the VBTF splay faults
over 30~yr of operation. Is it large enough to matter for seismic hazard?
\item Re-pointing economics: at what predicted decline rate does
re-completing or sidetracking become cheaper than drilling a new doublet?
This is a small but well-defined LCOH problem.
\end{itemize}

% ---------------------------------------------------------------------------
\subsection*{Which thread to pick?}

A pragmatic ranking, given a retreat-plus-followup timescale:

\begin{itemize}
\item \textbf{Thread 5 (geothermal vs.\ heat pumps)} is now the
highest-value thread in policy terms --- Wien Energie is actively
deploying both technologies in parallel and the implicit mix
question is genuinely open. A clean optimisation answer would be
publishable and useful.
\item \textbf{Thread 4 (storage coupling)} remains the most tractable
single-deliverable thread and the best aligned with Wien Energie's
stated open questions.
\item \textbf{Thread 7 (reservoir depletion)} is well-defined, has a
clean analytical core (thermal breakthrough scaling) and a strong
empirical baseline (Dogger 40-yr data); also feeds directly into
threads 2, 4, and 6.
\item \textbf{Thread 2 (seismicity)} is the most physics-rich and the
most pedagogically valuable --- rate-and-state friction, pore-pressure
diffusion, fault mechanics all in one place --- but the deliverable
is harder to wrap up cleanly.
\item \textbf{Thread 3 (lithium)} has the highest novelty ceiling
because the Vienna Basin Li data gap is real. A focused literature
review plus an order-of-magnitude bound is achievable and
publication-relevant.
\item \textbf{Thread 6 (CO$_2$ working fluid)} is the most speculative
but has just been operationally validated by Eavor at Geretsried
(Dec 2025), so the timing for a follow-up academic comparison is
ideal --- with the caveat that the thermodynamics are heavy and a
single retreat may not be enough.
\item \textbf{Thread 1 (well repurposing)} is the most concrete
because Aderklaa-96 is a live experiment with public data. Good if
the group prefers a single-case-study format.
\end{itemize}

Natural pairings:
\begin{itemize}
\item \textbf{Reservoir hydraulics pair:} threads 2, 4, and 7 all
share heat-and-pressure-diffusion modelling in the same aquifer
geometry. They can share a single underlying numerical model and
divide labour by which physics is foregrounded.
\item \textbf{Legacy-well pair:} threads 1, 3, and 6 all use the
existing OMV well inventory as their starting point (the wells
themselves; the brine chemistry inside them; the candidates for
sCO$_2$ retrofit respectively).
\item \textbf{Systems pair:} threads 4 and 5 share the same demand
curve and would benefit from a single shared optimisation framework.
\end{itemize}
A retreat that splits into three sub-groups along these three pairings
is a sensible structure, with the proviso that the underlying
shared-data work (see ``What is still missing'' below) needs to be
done jointly first.

% ---------------------------------------------------------------------------
\subsection*{What is still missing before this becomes a project plan}

\begin{itemize}
\item For threads 1 and 3 we still need access to operator data
(OMV production-geology archives, GeoSphere Austria / GBA reports).
Some of this may require an industry/agency contact rather than
just literature.
\item For thread 2, the GeoSphere Austria earthquake catalog should
be downloaded and cross-checked against the OMV oil and gas
production history. None of this requires permissions beyond what
is already public, but it requires data work nobody in the group
has yet done.
\item For thread 4, the Wien Energie demand curve at sub-daily
resolution would tighten the optimisation problem considerably ---
worth asking Wien Energie directly whether anonymised data are
available for academic use.
\item For thread 5, an hourly (or at least monthly) demand curve for
Vienna district heat plus an electricity price/carbon-intensity time
series for the same period are the two inputs that anchor the
optimisation; both should be assembled before retreat day 1.
\item For thread 6, locating any published or grey-literature sCO$_2$
geothermal results from the Eavor Geretsried commissioning would
sharpen the comparison considerably --- the very first commercial
operating data has only just been generated and is not yet in the
journal literature.
\item For thread 7, a Dogger-aquifer operating-data spreadsheet (Lopez
\emph{et al.}\ 2010 and the 2024 reanalysis) is the right empirical
calibration target; this is a literature-only ask but the
publication-quality figures need to be re-drawn from raw data.
\item All seven threads would benefit from a single afternoon's
spreadsheet collaboration to nail down shared numbers (Vienna heat
demand, doublet capacity, well inventory, electricity-price series,
Aderklaa Conglomerate properties) before splitting into sub-groups,
so the seven mini-projects rest on a common baseline.
\end{itemize}

% ---------------------------------------------------------------------------
\subsection*{References}

% IEEE-style bibliography. URLs use \texttt{} so they render in monospace
% even without the \texttt{url} package; if main.tex eventually loads
% \texttt{\textbackslash usepackage\{url\}} or \texttt{\textbackslash usepackage\{hyperref\}},
% \texttt{\textbackslash texttt} can be globally swapped for \texttt{\textbackslash url} for
% better line-breaking. Entries marked [verify] still need their author lists
% or full citations confirmed against the published versions.

\begin{thebibliography}{99}

\bibitem{ref1}
OMV AG, ``Climate-neutral district heating: beginning in 2028 --- Drilling
starts for first deep geothermal plant in Vienna,'' \emph{OMV press release},
Dec.~16, 2024. [Online]. Available:
\texttt{omv.com/en/media/press-releases/2024/241216-climate-neutral-district-heating-beginning-in-2028-drilling-starts-for-first-deep-geothermal-plant-in-vienna}.
[Accessed: May~19, 2026].

\bibitem{ref2}
ThinkGeoEnergy, ``Wien Energie to build first geothermal heating plant in
Vienna, Austria,'' 2023. [Online]. Available:
\texttt{thinkgeoenergy.com/wien-energie-to-build-first-geothermal-heating-plant-in-vienna-austria/}.
[Accessed: May~19, 2026].

\bibitem{ref3}
OMV AG, ``OMV starts two geothermal projects,'' \emph{OMV press release},
Oct.~3, 2022. [Online]. Available:
\texttt{omv.com/en/media/press-releases/2022/221003-omv-starts-two-geothermal-projects}.
[Accessed: May~19, 2026].

\bibitem{ref4}
ThinkGeoEnergy, ``Austrian OMV to test for geothermal at old natural gas
well,'' 2022. [Online]. Available:
\texttt{thinkgeoenergy.com/austrian-omv-to-test-for-geothermal-at-old-natural-gas-well/}.
[Accessed: May~19, 2026].

\bibitem{ref5}
``The Vienna Basin: Petroleum Systems, Storage and Geothermal Potential,''
edited volume / review chapter, 2024. ResearchGate Publication 383664592.
[Authors and full citation to verify.]

\bibitem{ref6}
E.~Hintersberger, K.~Decker, A.~Lomax, and C.~L\"uthgens, ``Implications from
palaeoseismological investigations at the Markgrafneusiedl Fault (Vienna
Basin, Austria) for seismic hazard assessment,'' \emph{Nat.\ Hazards Earth
Syst.\ Sci.}, vol.~18, pp.~531--553, 2018, doi: 10.5194/nhess-18-531-2018.

\bibitem{ref7}
W.~Lenhardt \emph{et al.}, ``The 2013 earthquake series in the Southern Vienna
Basin: Location and source,'' 2015. ResearchGate Publication 273983865.
[Full author list and venue to verify.]

\bibitem{ref8}
I.~Bianchi \emph{et al.}, ``The unusual seismic activity from 2021 to 2024 in
Eastern Austria: Insights from seismic sequences relocation and moment tensor
inversion,'' \emph{Bull.\ Seismol.\ Soc.\ Am.}, vol.~115, no.~1, p.~130,
2025. [Full author list to verify.]

\bibitem{ref9}
S.~Keil, J.~Wassermann, and T.~Megies, ``Estimation of ground motion due to
induced seismicity at a geothermal power plant near Munich, Germany, using
numerical simulations,'' \emph{Geothermics}, vol.~106, 2022.

\bibitem{ref10}
INSIDE Research Consortium, ``Investigating the relationship between
seismicity, deformation and deep geothermal exploitation in the Greater
Munich area,'' ResearchGate Publication 348232287. [Authors and year to
verify.]

\bibitem{ref11}
``Repurposing abandoned oil and gas wells for geothermal energy extraction:
A review of concepts, challenges, and opportunities,'' \emph{Appl.\ Energy},
Sep.\ 2025. ScienceDirect ID S0306261925015466. [Author list to verify.]

\bibitem{ref12}
X.~Hu, J.~Banks, L.~Wu, and W.~V.~Liu, ``Numerical modeling of a coaxial
borehole heat exchanger to exploit geothermal energy from abandoned petroleum
wells in Hinton, Alberta,'' \emph{Renewable Energy}, vol.~148,
pp.~1110--1123, 2020.

\bibitem{ref13}
N.~Leontidis \emph{et al.}, ``Controlling injection conditions of a deep
coaxial closed well heat exchanger to meet irregular heat demands: A field
case study in Belgium (Mol),'' \emph{Geothermal Energy}, vol.~13, art.~10,
2025, doi: 10.1186/s40517-025-00331-y.

\bibitem{ref14}
Electrive, ``Vulcan Energy mobilises EUR~2.2 billion for lithium extraction
in Germany,'' Dec.~3, 2025. [Online]. Available:
\texttt{electrive.com/2025/12/03/vulcan-energy-mobilises-e2-2-billion-for-lithium-extraction-in-germany/}.
[Accessed: May~19, 2026].

\bibitem{ref15}
Wikipedia contributors, ``Vulcan Energy Resources,'' \emph{Wikipedia, the Free
Encyclopedia}. [Online]. Available:
\texttt{en.wikipedia.org/wiki/Vulcan\_Energy\_Resources}.
[Accessed: May~19, 2026].

\bibitem{ref16}
International Lithium Association (ILiA), ``Direct lithium extraction
(DLE): An introduction,'' Tech.\ Rep., Jun.\ 2024. [Online]. Available:
\texttt{lithium.org/wp-content/uploads/2024/06/Direct-Lithium-Extraction-DLE-An-introduction-ILiA-June-2024-v.1-English-web.pdf}.
[Accessed: May~19, 2026].

\bibitem{ref17}
California Energy Commission, ``Pilot scale recovery of lithium from
geothermal brines,'' Tech.\ Rep.\ CEC-500-2024-020, Mar.\ 2024. [Online].
Available:
\texttt{energy.ca.gov/publications/2024/pilot-scale-recovery-lithium-geothermal-brines}.
[Accessed: May~19, 2026].

\bibitem{ref18}
S\&P Global Commodity Insights, ``Commodities 2026: Lithium carbonate surplus
to narrow; energy storage to drive growth,'' Jan.~9, 2026. [Online].
Available:
\texttt{spglobal.com/energy/en/news-research/latest-news/metals/010926-commodities-2026-lithium-carbonate-surplus-to-narrow-energy-storage-to-drive-growth}.
[Accessed: May~19, 2026].

\bibitem{ref19}
``Lithium enrichment processes in sedimentary formation waters,'' \emph{Chem.\
Geol.}, 2023. ScienceDirect ID S0009254123003261. [Author list to verify.]

\bibitem{ref20}
P.~Dobson, T.~McLing, N.~Spycher, P.~Fleuchaus, G.~Neupane, C.~Doughty,
Y.~Zhang, R.~Smith, T.~Atkinson, W.~Jin, P.~Blum, D.~Dinkelman, and
H.~Veldkamp, ``High-temperature aquifer thermal energy storage (HT-ATES)
projects in Germany and the Netherlands --- Review and lessons learned,''
\emph{Energies}, vol.~18, no.~23, art.~6292, 2025.

\bibitem{ref21}
D.~Zhou, K.~Li, H.~Gao, A.~Tatomir, M.~Sauter, and L.~Ganzer, ``Techno-economic
assessment of high-temperature aquifer thermal energy storage system,
insights from a study case in Burgwedel, Germany,'' \emph{Appl.\ Energy},
vol.~372, 2024. ScienceDirect ID S0306261924011668.

\bibitem{ref22}
R.~Vidal, S.~Olivella, M.~W.~Saaltink, and F.~Diaz-Maurin, ``Heat storage
efficiency, ground surface uplift and thermo-hydro-mechanical phenomena for
high-temperature aquifer thermal energy storage,'' \emph{Geothermal Energy},
2022, doi: 10.1186/s40517-022-00233-3.

\bibitem{ref23}
Green Energy Lab, ``ATES Vienna --- Aquifer thermal energy storage,'' project
page. [Online]. Available:
\texttt{greenenergylab.at/projects/ates-vienna}. [Accessed: May~19, 2026].

\bibitem{ref24}
TechXplore / AFP, ``Vast Vienna wastewater heat pumps showcase EU climate
drive,'' Mar.\ 2024. [Online]. Available:
\texttt{techxplore.com/news/2024-03-vast-vienna-wastewater-showcase-eu.html}.
[Accessed: May~20, 2026].

\bibitem{ref25}
Green Energy Lab / Wien Energie, ``Large heat pump in Vienna's Spittelau
district is now operational,'' 2026. [Online]. Available:
\texttt{greenenergylab.at/large-heat-pump-in-viennas-spittelau-district-now-operational-sustainable-heating-for-16000-households/}.
[Accessed: May~20, 2026].

\bibitem{ref26}
Heat Pumping Technologies Magazine, ``Roll-out of large-scale heat pumps as
a key factor for the German energy and heat transition,'' vol.~42, no.~2,
2024. [Online]. Available:
\texttt{heatpumpingtechnologies.org/heat-pumping-technologies-magazine-vol-42-no-2-2024-2/}.
[Accessed: May~20, 2026]. [Used here as a representative comparative-context
source; a direct academic LCOH comparison for large-DH applications appears
to be a literature gap and is itself a research question.]

\bibitem{ref27}
Eavor Technologies Inc., ``Eavor Technologies achieves first electricity
production at Geretsried site,'' \emph{Press release / GlobeNewswire},
Dec.~4, 2025. [Online]. Available:
\texttt{eavor.com/press-releases/eavor-technologies-achieves-first-electricity-production-at-geretsried-site/}.
[Accessed: May~20, 2026].

\bibitem{ref28}
\emph{Power Magazine}, ``Eavor's first-of-its-kind closed-loop geothermal
project produces grid power in Germany,'' Dec.\ 2025. [Online]. Available:
\texttt{powermag.com/eavors-first-of-its-kind-closed-loop-geothermal-project-produces-grid-power-in-germany/}.
[Accessed: May~20, 2026].

\bibitem{ref29}
B.~M.~Adams, D.~Vogler, T.~H.~Kuehn, J.~M.~Bielicki, N.~Garapati, and
M.~O.~Saar, ``Heat depletion in sedimentary basins and its effect on the
design and electric power output of CO$_2$ Plume Geothermal (CPG) systems,''
\emph{Renewable Energy}, vol.~172, pp.~1393--1403, 2021.

\bibitem{ref30}
S.~Lopez, V.~Hamm, M.~Le~Brun, L.~Schaper, F.~Boissier, C.~Cotiche, and
E.~Giuglaris, ``40 years of Dogger aquifer management in Ile-de-France,
Paris Basin, France,'' \emph{Geothermics}, vol.~39, no.~4, pp.~339--356,
2010.

\bibitem{ref31}
T.~Renaud, J.~Popineau, J.~O'Sullivan, and J.~Gasser Dorado, ``Novel
approach for modelling low enthalpy geothermal in deep sedimentary
aquifers: a case study of 40 years of production data in the Dogger
formation,'' \emph{Geothermal Energy}, vol.~12, art.~44, Dec.~19, 2024,
doi: 10.1186/s40517-024-00328-z.

\bibitem{ref32}
I.~Rahimzadeh Kivi, E.~Pujades, J.~Rutqvist, and V.~Vilarrasa,
``Cooling-induced reactivation of distant faults during long-term
geothermal energy production in hot sedimentary aquifers,'' \emph{Scientific
Reports}, vol.~12, art.~2065, 2022, doi: 10.1038/s41598-022-06067-0.

\end{thebibliography}

\paragraph{Note on the bibliography.}
This list was assembled from web searches in May~2026 and uses IEEE
formatting. Items marked ``[verify]'' had partial information returned by
the search (typically because the publisher page was paywalled or only
abstract metadata was indexed); their author lists and exact venue should
be confirmed against the published versions before circulation. The
publisher-confirmed academic references are [6], [9], [12], [13], [20],
[21], [22], [29], [30], [31], [32]; the OMV, Wien Energie / Green Energy
Lab, Tech Xplore, Eavor, Power Magazine, Heat Pumping Technologies,
Electrive, S\&P, ILiA, and CEC entries are primary-source online materials
and are reliable as written. Reference [26] is included as a
representative comparative-context source but is explicitly not an
academic LCOH comparison --- the absence of such a paper is itself one
of the open questions in thread~5.

\medskip
\noindent\emph{For a proper standalone bibliography file, these entries
can be migrated to BibTeX (\texttt{@article}, \texttt{@misc},
\texttt{@techreport} types) and rendered with \texttt{IEEEtran.bst}; the
inline \texttt{thebibliography} environment above produces the same visual
result without needing to run BibTeX as a separate compilation step.}
